\section{Comparison with the Original Results}

\subsection{Original Figure 4}
\begin{draftpoints}
  \item original Figure 4 places NDL with intrinsic replay slightly below CL with intrinsic replay in reconstruction error and shows qualitatively more stable reconstructed digits \parencite{draelos2017neurogenesis}
  \item original plot is treated as directional evidence because the underlying seed-level numerical table and confidence intervals are not available in the preserved paper source
  \item replication therefore tests whether the direction of the matched comparison is reproduced, rather than claiming an exact difference from values estimated visually from the published graphic
\end{draftpoints}

\begin{figure}[htbp]
\centering
\includegraphics[width=0.70\textwidth]{original_figure4.png}
\caption{Original paper Figure 4, reproduced for directional comparison. Panels compare CL, NDL, CL+IR, and NDL+IR and show the reported final reconstruction advantage of NDL+IR \parencite{draelos2017neurogenesis}.}
\label{fig:original-figure4}
\end{figure}

\subsection{Claim-by-claim verdict}
\begin{draftpoints}
  \item Mechanism implementation is substantially reproduced: reconstruction-triggered layer growth, protected mature features, reduced-rate decoder updates, replay stabilization, and next-layer propagation are present
  \item Architectural resemblance can be reproduced, but the closest paper-like endpoint is often explained by configured allowances rather than successful quota-based stopping
  \item central numerical ordering is not reproduced: NDL+IR is 2.70 times worse than CL+IR in macro MSE and 2.06 times worse in positive forgetting
  \item reported SD-19 advantage is not reproduced in the three-seed feasibility comparison, where matched CL is better in reconstruction, forgetting, updates, and runtime
  \item expected constant-factor overhead is not supported by the endpoint-matched resource benchmark, especially for intrinsic replay
\end{draftpoints}

\begin{table}[htbp]
\centering
\caption{Original claims versus replication observations. ``Directional'' indicates that the original paper does not expose exact seed-level values for a numerical equality test.}
\label{tab:original-versus-replication}
\small
\begin{tabularx}{\textwidth}{p{3.0cm}YYp{2.4cm}}
\toprule
Claim & Original conclusion & Replication observation & Status \\
\midrule
NDL+IR beats CL+IR & Slight reconstruction advantage & 2.70$\times$ higher macro MSE & Not replicated \\
NDL+IR preserves old classes & Better qualitative stability & 2.06$\times$ positive forgetting & Not replicated \\
Growth discovers useful size & Capacity added as needed & Many loops end at allowances with unresolved outliers & Not demonstrated \\
No-replay growth improves continual learning & Growth protects representations & Forgetting improves, total MSE worsens 1.91$\times$ & Partial trade-off \\
SD-19 supports growth & NDL+IR below CL+IR error & NDL about 1.85$\times$ macro MSE & Not replicated \\
Overhead remains moderate & Constant-factor argument & Up to 6.88$\times$ runtime in controlled benchmark & Not supported \\
\bottomrule
\end{tabularx}
\end{table}

\begin{figure}[htbp]
\centering
\includegraphics[width=0.94\textwidth]{cross_dataset_effects.png}
\caption{Cross-dataset ratios for matched NDL and CL comparisons. Ratios above one indicate worse NDL reconstruction, forgetting, or update cost. Source: \artifact{cross_dataset_ratios.csv}.}
\label{fig:cross-dataset-effects}
\end{figure}

\subsection{Interpretive boundary}
\begin{draftpoints}
  \item evidence rejects a numerical replication claim for this implementation but does not establish that neuron growth can never improve continual learning
  \item publication's underspecified optimizer, activation, threshold, quota, phase schedule, and node-count policy leave a family of plausible implementations rather than one exactly recoverable experiment
  \item negative matched comparisons show that reproducing the stated mechanism and approximate endpoint shape is insufficient to reproduce the reported performance advantage
\end{draftpoints}
